% ----------------------------------------------------------------------
%  2.1.3. Комплексы точечных дефектов
% ----------------------------------------------------------------------
\subsection{Комплексы точечных дефектов}\label{ssec:complexes}

\begin{orientir}
Точечные дефекты могут объединяться: две вакансии~— в дивакансию, три~—
в тривакансию и т.\,д. Это энергетически выгодно. Здесь мы выясним,
сколько таких комплексов в равновесии, как они устроены и насколько они
подвижны. Главный практический вывод: комплексы из трёх и более вакансий
малоподвижны и служат зародышами пор.
\end{orientir}

\subsubsection{Энергия связи и равновесная концентрация комплексов}

Почему две вакансии «хотят» быть рядом? Вспомним модель разрыва связей.
Чтобы создать изолированную вакансию, нужно разорвать $z$ связей. Если же
мы вынимаем атом, \emph{соседний} с уже существующей вакансией, одна из
его связей уже разорвана~— рвать приходится на одну связь меньше. Значит,
пара соседних вакансий стоит меньше энергии, чем две отдельные.

\begin{opredelenie}
\term{Дивакансия}~— пара вакансий в соседних узлах. \term{Энергия связи}
вакансий в дивакансии $B_2$~— выигрыш энергии при объединении двух
изолированных вакансий в дивакансию:
\begin{equation}\label{eq:E2v}
E_{2v}^{\mathrm F}=2\EvF-B_2 .
\end{equation}
\end{opredelenie}

\paragraph{Концентрация дивакансий}
Повторим рассуждение п.~\ref{sssec:eq-conc}, но теперь «местами» для
дефекта служат не узлы, а пары соседних узлов.
\begin{shagi}
\item У каждого из $N$ узлов $z$ соседей, и каждая пара учтена дважды;
значит, пар соседних узлов $zN/2$.
\item Число способов разместить $n_{2v}$ дивакансий по этим парам:
\[
W=\frac{(zN/2)!}{\bigl(zN/2-n_{2v}\bigr)!\;n_{2v}!}\,.
\]
\item Минимизация $\Delta G=n_{2v}E_{2v}^{\mathrm F}-kT\ln W$, как и раньше,
даёт
\[
C_{2v}=\frac{n_{2v}}{N}=\frac{z}{2}\Ex{-\frac{E_{2v}^{\mathrm F}}{kT}} .
\]
\item Подставим~\eqref{eq:E2v} и учтём, что $\exp(-\EvF/kT)=C_{1v}$
(концентрация одиночных вакансий, моновакансий):
\end{shagi}
\begin{keyeq}\label{eq:C2v}
C_{2v}=\frac{z}{2}\,C_{1v}^2\Ex{\frac{B_2}{kT}},\qquad
\text{для ГЦК ($z=12$):}\quad \frac{C_{2v}}{C_{1v}}=6\,C_{1v}\Ex{\frac{B_2}{kT}} .
\end{keyeq}
Множитель $6=z/2$~— число независимых ориентаций дивакансии в ГЦК-решётке.

\begin{smysl}
$C_{2v}\propto C_{1v}^2$: дивакансия~— это «встреча» двух вакансий, и
вероятность её пропорциональна квадрату их концентрации. Множитель
$\exp(B_2/kT)$ показывает, во сколько раз встреча выгоднее случайной
благодаря притяжению.
\end{smysl}

\begin{primer}[Доля дивакансий у точки плавления]
Медь, $T=\Tpl=1358$~К ($kT=0{,}117\eV$), $C_{1v}\approx10^{-4}$, $B_2\approx
0{,}2\eV$:
\[
\frac{C_{2v}}{C_{1v}}=6\cdot10^{-4}\cdot e^{0{,}2/0{,}117}\approx
6\cdot10^{-4}\cdot5{,}5\approx 3\cdot10^{-3}.
\]
Даже у точки плавления дивакансий меньше процента от числа моновакансий.
\end{primer}

\begin{glavnoe}
В тепловом равновесии в чистых металлах почти все вакансии~—
одиночные: даже у точки плавления моновакансии составляют более 90\,\%.
Роль дивакансий и более крупных комплексов резко возрастает, когда
кристалл \emph{пересыщен} вакансиями (после закалки, облучения): при
большом $C_{1v}$ квадратичная зависимость~\eqref{eq:C2v} даёт много пар.
\end{glavnoe}

Те же рассуждения применимы к любому комплексу (тривакансии, тетравакансии
и т.\,д.): концентрация комплекса из $m$ вакансий пропорциональна
$C_{1v}^m\exp(B_m/kT)$.

\subsubsection{Строение комплексов}

Детальные расчёты показывают, что комплексы вакансий и межузельных атомов
устойчивы, причём их строение бывает неожиданным. Рассмотрим три
возможные конфигурации тривакансии в ГЦК-меди (рис.~\ref{fig:trivac}).

\begin{figure}[ht]
\centering
\begin{tikzpicture}[cubeproj, scale=2.7]
  % (а) линейная
  \begin{scope}[shift={(0,0,0)}]
    \foreach \t in {0,1,2}{\node[vac=4.4mm] at ({0.5*\t},{0.5*\t},0) {};}
    \draw[atomLine!50, dashed] (0,0,0) -- (1,1,0);
    \node[lbl] at (0.5,0,-0.4) {а};
  \end{scope}
  % (б) плоская
  \begin{scope}[shift={(1.8,0,0)}]
    \coordinate (p1) at (0,0,0); \coordinate (p2) at (0.5,0.5,0);
    \coordinate (p3) at (0.5,0,0.5); \coordinate (p4) at (0,0.5,0.5);
    \draw[atomLine!70, line width=0.6pt] (p1)--(p2)--(p3)--cycle (p1)--(p4)--(p2) (p3)--(p4);
    \node[vac=4.4mm] at (p1) {}; \node[vac=4.4mm] at (p2) {}; \node[vac=4.4mm] at (p3) {};
    \node[atomB=4.4mm] at (p4) {}; \node[slbl, above right] at (p4) {A};
    \node[lbl] at (0.25,0,-0.4) {б};
  \end{scope}
  % (в) тетраэдрическая
  \begin{scope}[shift={(3.4,0,0)}]
    \coordinate (p1) at (0,0,0); \coordinate (p2) at (0.5,0.5,0);
    \coordinate (p3) at (0.5,0,0.5); \coordinate (p4) at (0,0.5,0.5);
    \coordinate (c) at (0.25,0.25,0.25);
    \draw[atomLine!70, line width=0.6pt] (p1)--(p2)--(p3)--cycle (p1)--(p4)--(p2) (p3)--(p4);
    \foreach \p in {p1,p2,p3,p4}{\node[vac=4.4mm] at (\p) {};}
    \draw[arr, fmRed, line width=0.5pt] (p4) -- ($(p4)!0.75!(c)$);
    \node[atomB=4.4mm] at (c) {}; \node[lbl] at (0.25,0,-0.4) {в};
  \end{scope}
\end{tikzpicture}
\caption{Три конфигурации тривакансии в ГЦК-решётке (пунктир~— пустые узлы):
\textit{а}~— линейная; \textit{б}~— плоская: три вакансии~— ближайшие
соседи; \textit{в}~— тетраэдрическая: атом~A сместился в центр
тетраэдра, образованного четырьмя узлами}
\label{fig:trivac}
\end{figure}

\begin{enumerate}[label=\asbuk*)]
\item \emph{Линейная} (рис.~\ref{fig:trivac},\,а): крайние вакансии не
являются ближайшими соседями, выигрыш от связи мал~— конфигурация
наименее устойчива.
\item \emph{Плоская} (рис.~\ref{fig:trivac},\,б): все три вакансии~—
ближайшие соседи. Казалось бы, это и есть самая устойчивая форма.
\item \emph{Тетраэдрическая} (рис.~\ref{fig:trivac},\,в): расчёт
показывает, что атом~А, общий сосед трёх вакансий, «проваливается» в
центр тетраэдра. Получается необычный объект: три вакансии «размазаны»
поровну по четырём вершинам тетраэдра, а в его центре сидит смещённый
атом. Эта релаксация даёт большой вклад в энергию связи, и именно такая
тривакансия устойчива.
\end{enumerate}

\begin{glavnoe}
Устойчивость крупных комплексов определяется не только числом связей
между вакансиями, но и \emph{сильной релаксацией} отдельных атомов,
которые смещаются в несимметричные положения. Так стабилизируются,
например, октаэдрические формы тетра- и пентавакансий: в них один или
два атома сильно смещены к центру комплекса.
\end{glavnoe}

Для сдвоенных межузельных атомов расчёт находит три устойчивые
конфигурации с очень близкими энергиями связи (около $0{,}5\eV$), так что
указать преобладающую нельзя.

\subsubsection{Подвижность комплексов}

\paragraph{Дивакансия подвижнее вакансии}
Зейтц впервые заметил, что в плотноупакованных металлах дивакансия должна
быть \emph{более} подвижной, чем одиночная вакансия. Главный вклад в
энергию миграции даёт отталкивание атомов, «запирающих» путь
перескакивающему атому (рис.~\ref{fig:vac-jump},\,б). Если один из
запирающих атомов отсутствует (а это и есть вторая вакансия), то,
во-первых, исчезает одно отталкивание, а во-вторых, перескакивающий атом
может пройти по пути, более удалённому от оставшихся. Расчёт для меди даёт
энергии активации миграции $\approx0{,}97\eV$ для моновакансии и лишь
$\approx0{,}35\eV$ для дивакансии.

\paragraph{Тривакансия малоподвижна}
Чтобы тетраэдрическая тривакансия переместилась, центральный атом должен
вернуться в узел, а другой атом~— сместиться в центр соседнего
тетраэдра; комплекс при этом частично «разбирается». Энергия активации
такого процесса около $1{,}9\eV$~— примерно в полтора раза больше, чем у
моновакансии, и намного больше, чем у дивакансии. Основной вклад в неё
дают как раз большие смещения атомов при перестройке.

\begin{glavnoe}
Тривакансия~— наименьший комплекс вакансий, который одновременно
устойчив и малоподвижен. Поэтому её считают \emph{зародышем
минимального размера} при образовании вакансионных пор: к неподвижной
тривакансии подходят новые вакансии, и комплекс растёт. Крупные комплексы
ещё менее подвижны.
\end{glavnoe}

Сдвоенные межузельные атомы также менее подвижны, чем одиночные.

\begin{itogi}
\begin{itemize}
\item Вакансии притягиваются: $E_{2v}^{\mathrm F}=2\EvF-B_2$.
\item $C_{2v}=\frac{z}{2}C_{1v}^2\exp(B_2/kT)$; в равновесии почти все
вакансии одиночные, но при пересыщении доля комплексов резко растёт.
\item Устойчивая тривакансия в ГЦК~— тетраэдрическая с атомом в центре;
сильная релаксация отдельных атомов стабилизирует комплексы.
\item Подвижность: дивакансия подвижнее моновакансии; тривакансия и
более крупные комплексы малоподвижны и служат зародышами пор.
\end{itemize}
\end{itogi}

\begin{proverka}
\item Почему энергия образования дивакансии меньше суммы энергий двух
вакансий?
\item Почему концентрация дивакансий пропорциональна квадрату
концентрации моновакансий?
\item Почему дивакансия подвижнее одиночной вакансии, а тривакансия~—
наоборот, малоподвижна?
\end{proverka}
